\clearpage
\thispagestyle{plain}
\begin{center}
\textbf{Contents}
\end{center}

\vspace{0.3in}

\begin{tabular}{@{}p{0.75in}p{4.35in}r@{}}
1 & Introduction & 1\\[0.3em]
2 & Implementing the Complete \(\lambda\)-calculus & 3\\
2.1 & The Pure Untyped \(\lambda\)-Calculus & 3\\
2.1.1 & Syntax & 3\\
2.1.2 & Semantics & 4\\
2.1.3 & Normal Form & 5\\
2.2 & Implementation Issues & 5\\
2.2.1 & Representation of Expressions & 5\\
2.2.2 & Incremental Reduction & 6\\
2.2.3 & Order of Reduction & 7\\
2.2.4 & Substitution & 7\\
2.3 & The Theory of Head Order Reduction & 11\\
2.4 & Deriving an Architecture & 12\\
2.4.1 & DeBruijn Indices & 13\\
2.5 & Alternative Derivations & 14\\[0.3em]
3 & The Operational Semantics of THOR & 15\\
3.1 & The Syntax of Thor & 16\\
3.2 & Notation & 17\\
3.2.1 & Configurations & 17\\
3.2.2 & Reduction Rules & 17\\
3.2.3 & Miscellaneous & 18\\
3.3 & The Pure \(\lambda\)-calculus & 19\\
3.4 & Adding in Data Objects & 20\\
3.4.1 & Constants & 20\\
3.4.2 & Primitive Operations on Constants & 21\\
3.4.3 & Structures & 22\\
3.5 & Logical Primitives & 23\\
3.6 & The Conditional Expression & 25\\
3.7 & Recursion & 26\\
3.7.1 & The \(Y\) Operator & 26\\
3.7.2 & Mutual Recursion with LETREC & 28\\
3.7.3 & Recursion via Symbol Definitions & 30\\
3.7.4 & Controlling Recursion & 30\\
3.8 & The \(\alpha\)-Equality Predicate & 31\\
\end{tabular}

\clearpage
\thispagestyle{plain}
\begin{tabular}{@{}p{0.75in}p{4.35in}r@{}}
4 & The RED2 Architecture & 34\\
4.1 & The Simple Stack Machine (SSM) & 34\\
4.1.1 & Representing Expressions & 34\\
4.1.2 & The Abstract Machine & 37\\
4.1.3 & Correctness of the Simple Stack Machine & 39\\
4.2 & Mapping the SSM onto Hardware: \(\mu\)RED & 45\\
4.2.1 & The Instruction Set & 45\\
4.2.2 & The Memory Model & 46\\
4.2.3 & Registers & 47\\
4.2.4 & Register Transfer Notation & 47\\
4.2.5 & The Environment & 47\\
4.2.6 & The Execution Model & 49\\
4.2.7 & Implementing the Control Unit & 50\\
4.2.8 & Optimizing Single Instruction Arguments & 56\\
4.3 & RED2 & 58\\
4.3.1 & The Head Flag & 58\\
4.3.2 & Adding Basic Data Types & 59\\
4.3.3 & Symbolic Constants & 59\\
4.3.4 & Built-In Primitives & 61\\
4.3.5 & Implementing Lazy Structures & 66\\
4.3.6 & Local Mutual Recursion & 67\\[0.3em]
5 & Analysis and Enhancement of RED2 & 72\\
5.1 & Comparing RED2 with the Three Instruction Machine & 72\\
5.2 & Sharing & 72\\
5.2.1 & Why RED2 Can't Always Share & 74\\
5.2.2 & What RED2 Can Share & 76\\
5.2.3 & Extending RED2 to Increase Sharing & 77\\
5.2.4 & Detecting Closed Expressions & 79\\
5.3 & Towards Implementing \(\eta\)-Reduction & 79\\
5.4 & Enhanced Compilation and von-Neumann Plus Architectures & 82\\
\end{tabular}

\clearpage
\thispagestyle{plain}
\begin{tabular}{@{}p{0.75in}p{4.35in}r@{}}
6 & Implementing Horn Clause Logic on RED2 & 84\\
6.1 & Logic Programming & 84\\
6.2 & Integrating THOR and Logic Programming & 85\\
6.2.1 & A Uniform Notation & 85\\
6.2.2 & Unification & 85\\
6.2.3 & Search & 86\\
6.2.4 & Interfacing Functions and Relations & 87\\
6.3 & The Unification Algorithm & 87\\
6.4 & Implementing Unification on RED2 & 89\\
6.4.1 & Indirection Pointers & 90\\
6.4.2 & Backtracking and the Reset List & 90\\
6.4.3 & Unification Memory & 90\\
6.4.4 & EP's and Delayed Variable Binding & 91\\
6.5 & Managing Explicit Quantification & 92\\
6.5.1 & Removing LAMBDA's Bound Through Unification & 92\\
6.5.2 & Extending the Scope of Quantifiers & 93\\
6.6 & Adapting Unification to Incremental Reduction & 93\\
6.7 & Negation & 94\\
6.8 & Generating All Solutions of a Goal & 94\\
6.9 & Discussion & 96\\[0.3em]
A & Benchmark Programs & 97\\
A.1 & The SINE Program & 97\\
A.2 & The GAME Program & 98\\
\end{tabular}
